\section{Science and Evaluation}
\sectionkicker{DISCOVERY \& BENCHMARKS}
\begin{wideflow}
{\Large\bfseries 確かめの道具立て――酵素と二つの物差し}\par
\smallskip
{\color{SurveyMuted}9月21～23日の科学と評価の三件を、未検証の線引きとともに示す。}\par
\end{wideflow}

9月23日、Anthropicは生き物の研究所の構えと最初の実りを示した\cite{enzymeart}。湾岸の濡れ場（BSL-1とBSL-2、腕は人の研究者）を添え、Claudeの使い手およそ950が21時間、2億1000万トークンをかけてDNAの海を探り、逆転写酵素の変わり種の傍らにCRISPRを思わせる繰り返しの列を見つけたとされる。列と添えの種からなるARTと名づけ、列が短いRNAとして読まれることは確かめたが、働きは分からないと明言する。書きぶりは査読前の草稿であり、中身の正しさは本号では量らない。最初の問いと濡れ仕事は人の手で、探し回りが使い手の領分という線引きを崩さない。CRISPRの開拓者の一人（Feng Zhang）が草稿を読み、調べるに値すると述べたことも添えられるが、追試ではない。

同じ23日、OpenAIは心の寄り添いの物差しMentalHealthBenchを開いた\cite{mhb}。22か国80人超の免許持ち（19言語、20近い分科）が作り、1,215の作り話のやり取りに、軽い日常から切迫まで、大人・若年・介護・臨床の顔ぶれを並べる。一つの返しに-10から+10の目盛りの規準をつけ、三人以上で読み、二人の一致と第三者の否定なしで残す流儀とされる。採点の手はGPT-5.6 Sol（自家の品で自家の物差しを採点する巡り）に委ねられ、点の比べは本号では真としない。利用者の声（44人）は調子と次の一手を重んじ、専門家は文脈の聞き取りを重んじる向きの違いも添えられる。

9月21日（改め22日）にはmem0から使い手の記憶の物差しDolphinBenchが出た\cite{dolphin}。知識仕事の三つの顔に各50万トークンほどの来歴を持たせ、200題ずつ、来歴ありで解けて来歴なしで解けないことをもって題の要否を確かめ、正しさに費用と待ち時間を添えて比べる建て付けとされる。6枚の草稿の表紙（arXiv）までしか読んでおらず、題の質や土台の数は量らない。公開の場の26件にはこの件の扱える投稿がなく、草稿の権威だけで立つ。

\begin{claimboundary}[資料と評価の境界]
ARTの働きや遺伝子操作への用立ては分からず、科学の万能の証しには使わない。物差しの点比べは真としない。DolphinBenchの題の質や土台の数は本号の検証範囲外である。
\end{claimboundary}
